\appendices{}

\section{Relative Motion in a Circular Orbit with No Drag} \label{app:cwh-equations}
This appendix references \cite{clohessy1960terminal} to describe the motion of a chaser in close proximity to a target in a circular orbit with no drag.
The translational dynamics of the chaser are described using the Clohessy-Wiltshire-Hill (CWH) Equations as follows:
\begin{subequations}
\begin{align}
    \frac{d^2x}{dt^2}&=2n\frac{dy}{dt}\\
    \frac{d^2y}{dt^2}&=-2n\frac{dx}{dt}+3n^2y\\
    \frac{d^2z}{dt^2}&=-n^2z,
\end{align}
\end{subequations}
where the x-axis points radially into the Earth (opposite to the radius vector of the target spacecraft), the z-axis is along the angular momentum vector of the target spacecraft, and the x-axis completes the right-handed system.
Here,
\begin{align}
    n&=\sqrt{\frac{\mu}{a}},
\end{align}
where $\mu=GM_E$ is the standard gravitational parameter for the Earth as the central body (geocentric gravitational constant) and $a$ is the radius of the target's orbit.
This can be rewritten as a linear time-invariant (LTI) system:
\begin{equation}
    \dot{\boldsymbol{r}}(t)=A\boldsymbol{r}(t)+B\boldsymbol{u}(t)
\end{equation}
where
\begin{subequations}
\begin{align}
    \boldsymbol{r}(t)&=\begin{bmatrix}
    x(t)&y(t)&z(t)&\dot{x}(t)&\dot{y}(t)&\dot{z}(t)
    \end{bmatrix}^T,\\
    \boldsymbol{u}(t)&=\begin{bmatrix}
    u_x(t)&u_y(t)&u_z(t)
    \end{bmatrix}^T,\\
    A&=\begin{bmatrix}
        0_{3\times3}&I_{3\times3}\\
        n^2A_{\dot{r}r}&nA_{\dot{r}\dot{r}}
    \end{bmatrix},\\
    B&=\begin{bmatrix}
        0_{3\times3}\\
        I_{3\times3}
    \end{bmatrix}
\end{align}
\end{subequations}
such that:
\begin{subequations}
\begin{align}
    A_{\dot{r}r}&=\begin{bmatrix}
        0&0&0\\
        0&3&0\\
        0&0&-1
    \end{bmatrix},\\
    A_{\dot{r}\dot{r}}&=\begin{bmatrix}
        0&2&0\\
        -2&0&0\\
        0&0&0
    \end{bmatrix}.
\end{align}
\end{subequations}

\section{Relative Motion in an Elliptical Orbit with Quadratic Drag} \label{app:ellp-undrag}
This appendix references \cite{carter2002clohessy} to describe the motion of a chaser with respect to the local-vertical-local-horizontal frame of reference of the target in an elliptical orbit with quadratic drag (considering the possibility of unequal drag coefficients between the chaser and the target). The dynamics of the chaser are given as:
\begin{equation}
    \boldsymbol{r}'(\theta)=A(\theta,\theta_0)\boldsymbol{r}(\theta)+B\boldsymbol{u}(\theta)+\boldsymbol{d}(\theta, \theta_0)
\end{equation}
where
\begin{subequations}
\begin{align}
    \boldsymbol{r}(\theta)&=\begin{bmatrix}
    x(\theta),y(\theta),z(\theta),x'(\theta),y'(\theta),z'(\theta)
    \end{bmatrix}^T,\\
    \boldsymbol{u}(\theta)&=\begin{bmatrix}
    u_x(\theta),u_y(\theta),u_z(\theta)
    \end{bmatrix}^T\\
    A(\theta, \theta_0)&=\begin{bmatrix}
        0_{3\times3}&I_{3\times3}\\
        A_{r'r}(\theta, \theta_0)&A_{r'r'}(\theta, \theta_0)
    \end{bmatrix},\\
    B&=\begin{bmatrix}
        0_{3\times3}\\
        I_{3\times3}
    \end{bmatrix},\text{ and}\\
    d(\theta, \theta_0)&=\begin{bmatrix}
        0_{3\times1}\\
        R(\theta,\theta_0)d_i(\theta,\theta_0)\\
        -R'(\theta,\theta_0)d_i(\theta,\theta_0)\\
        0
    \end{bmatrix}
\end{align}
\end{subequations}
such that:
\begin{subequations}
\begin{align}
    A_{r'r}&=\begin{bmatrix}
        -\gamma\frac{R'}{R}&\gamma&0\\
        -\gamma&-\gamma\frac{R'}{R} + \frac{3Re^{2\alpha\theta}\mu}{h^2}&0\\
        0&0&-1-\gamma\frac{R'}{R}
    \end{bmatrix}\\
    A_{r'r'}&=\begin{bmatrix}
        -\gamma&2&0\\
        -2&-\gamma&0\\
        0&0&\gamma
    \end{bmatrix},
\end{align}
\end{subequations}
and the new quantities dependent on eccentricity and drag coefficients are defined as
\begin{subequations}
\begin{align}
    \gamma&= \beta - \alpha\\
    d_i(\theta,\theta_0) &= \frac{\gamma R(\theta,\theta_0)e^{\alpha\theta}}{\sqrt{h}}\\
    R(\theta,\theta_0) &= \frac{h^2(1+4\alpha^2)}{\mu} \frac{1}{e^{2\alpha\theta} + {\mathsf{e}}\cos{(\theta - \theta_0)}}\\
    R'(\theta,\theta_0)&=\frac{dR}{d\theta}(\theta,\theta_0)
\end{align}
\end{subequations}
where ${\mathsf{e}}$ is the eccentricity of the orbit of the target, $\theta$ is the true anomaly of the orbit of the target, $\alpha$ is the drag coefficient of the target, and $\beta$ is the drag coefficient of the chaser.
We point the reader to observe the equivalence of this system to the CWH equations provided in Appendix \ref{app:cwh-equations} when the eccentricity of the orbit and the drag coefficients of both the target and the chaser are all set to zero.\\
We draw a parallel between the drag equation provided in the paper:
\begin{equation}
    \boldsymbol{a}_{D}=-\frac{\alpha}{R}|\dot{\boldsymbol{R}}|\dot{\boldsymbol{R}}
\end{equation}
where $\boldsymbol{R}$ is the position of the satellite from the center of the Earth and $R=|\boldsymbol{R}|$, and the standard drag equation \cite{batchelor2000an}
\begin{equation}
    F_{D}=\frac{1}{2}\rho v^{2}C_{D}A
\end{equation}
where $F_{D}$ is the drag force, $\rho$ is the density of air, $v$ is the speed of the chaser, $A$ is the cross sectional area, and $C_{D}$ is the dimensionless drag coefficient of the chaser to deduce that $\alpha$ and $\beta$ lie in the range of $10^{-7}$ units.
This is used in preparing the setup for the results we present in Section \ref{sec:results}

\section{The Minimum Enclosing Ellipsoid of the Target}
\label{app:mee}
We consider the minimum enclosing ellipsoid (MEE) \cite{kallberg2019minimum} of the target to be a pseudo-target of radii $a$, $b$, and $c$.
Let T be a set of points on the target that minimally defines the complete boundary of the target.
The minimum enclosing ellipsoid of the target can be defined as:
\begin{align}
    \frac{x^2}{a^2}+\frac{y^2}{b^2}+\frac{z^2}{c^2}&=1
\end{align}
where $a$, $b$, and $c$ can be computed using the following minimization problem:
\begin{equation}
\begin{aligned}
    \min_{a, b, c}\quad &V=\frac{4}{3}\pi abc\\
    \textrm{s.t. } & a,b,c>0\\
    & \frac{p_x^2}{a^2}+\frac{p_y^2}{b^2}+\frac{p_z^2}{c^2}<1\quad\forall\ p\in \text{T}
    \label{eqn:mee}
\end{aligned}
\end{equation}
